
\begin{obeactivity}
\textbf{Aktivitas 7.1: Menghitung Waktu Pencarian Kunci}
1. Hitung waktu terburuk dan waktu rata-rata \textit{brute force} DES (56 bit)
   jika satu komputer dapat mencoba $10^9$ kunci per detik. Gunakan $2^{56}$
   kemungkinan kunci dan waktu rata-rata $2^{55}$ percobaan.
2. Bandingkan hasilnya dengan waktu yang diperlukan untuk mencoba seluruh kunci
   Triple DES tiga kunci, yaitu $2^{168}$ kemungkinan, pada kecepatan yang sama.
3. Diskusikan mengapa mengulang DES dua kali dengan dua kunci (Double DES) tidak
   menambah keamanan sebesar yang diharapkan.
\end{obeactivity}

\begin{obeactivity}
\textbf{Aktivitas 7.2: Mengerjakan Satu Putaran DES di Atas Kertas}
Kerjakan langkah-langkah berikut tanpa program, memakai tabel pada
Tabel~\ref{tab:ip-des}, Tabel~\ref{tab:ep-des}, Tabel~\ref{tab:sbox-des}, dan
Tabel~\ref{tab:pc-des}. Gunakan plainteks $\code{0123456789ABCDEF}$ dan kunci
$\code{133457799BBCDFF1}$.

1. Baca tabel PC-1 pada kunci untuk memperoleh $C_0$ dan $D_0$. Pastikan bahwa
   delapan bit paritas benar-benar terbuang.
2. Geser $C_0$ dan $D_0$ satu bit ke kiri, lalu baca PC-2 untuk memperoleh $K_1$.
   Bandingkan hasil Anda dengan $K_1$ pada Contoh~\ref{ex:des-subkunci}.
3. Baca tabel IP pada plainteks untuk memperoleh $L_0$ dan $R_0$.
4. Baca tabel $E$ pada $R_0$, lalu XOR hasilnya dengan $K_1$.
5. Potong hasilnya menjadi delapan kelompok enam bit dan substitusikan melalui
   $S_1$ sampai $S_8$.
6. Baca tabel $P$ pada hasil kotak-S, lalu XOR dengan $L_0$ untuk memperoleh
   $R_1$.
7. Bandingkan seluruh hasil kerja Anda dengan tabel jejak putaran pada
   Contoh~\ref{ex:des-putaran-pertama}.
\end{obeactivity}

\begin{obeactivity}
\textbf{Aktivitas 7.3: Memeriksa Vektor Uji dan Kunci Lemah dengan Program}
Berkas \texttt{code/python/des\_uji.py} berisi implementasi DES lengkap yang
memuat seluruh tabel bit pada bab ini. Jalankan program itu
(\texttt{python des\_uji.py}) dan perhatikan keluarannya:

1. Periksa bahwa program menghasilkan vektor uji resmi
   $\code{85E813540F0AB405}$ dan bahwa dekripsi mengembalikan plainteks semula.
2. Bandingkan jejak putaran $(L_i, R_i)$ yang dicetak program dengan tabel pada
   Contoh~\ref{ex:des-putaran-pertama}. Keduanya harus identik.
3. Amati bagian kunci lemah: keempat kunci pada daftar itu menghasilkan 16
   subkunci yang sama. Jelaskan mengapa hal itu terjadi dengan menelusuri
   $C_0$ dan $D_0$ untuk kunci $\code{0101010101010101}$.
4. Amati bagian efek longsoran dan catat berapa bit cipherteks yang berubah
   ketika satu bit plainteks dibalik.
\end{obeactivity}
